\documentclass[11pt]{article}
\usepackage[letterpaper,margin=0.85in]{geometry}
\usepackage{mathpazo}
\usepackage[T1]{fontenc}
\usepackage[utf8]{inputenc}
\usepackage{microtype}
\usepackage{enumitem}
\usepackage{booktabs}
\usepackage{tabularx}
\usepackage{array}
\usepackage{fancyhdr}
\usepackage[hidelinks]{hyperref}
\usepackage{xurl}
\setlength{\parindent}{1.25em}
\setlength{\parskip}{0.2em}
\setlist{nosep,leftmargin=1.5em}
\renewcommand{\arraystretch}{1.18}
\newcolumntype{Y}{>{\raggedright\arraybackslash}X}
\pagestyle{fancy}
\fancyhf{}
\lhead{API 209 Math Camp 2026}
\rhead{Lab 4: Quarto policy brief}
\cfoot{\thepage}
\setlength{\headheight}{14pt}
\newcommand{\file}[1]{\nolinkurl{#1}}
\begin{document}

\begin{center}
{\small WEDNESDAY, AUGUST 26 \quad 4:00--5:00 P.M.}\par
\vspace{0.6em}
{\LARGE\textbf{Lab 4: From a policy question to a Quarto brief}}\par
\vspace{0.4em}
{\large One table, one figure, one finding, and one limitation}
\end{center}

\section*{Purpose}

Choose one of four international-development questions and complete its
Quarto starter. Your brief should let a reader identify the sample, inspect
the calculation, and distinguish a descriptive result from an unsupported
claim. You will write the first interpretation before asking Codex to review
the work.

\section*{Choose one track}

\begin{tabularx}{\textwidth}{@{}p{0.15\textwidth}Y p{0.22\textwidth}@{}}
\toprule
Track & Question & Starter \\
\midrule
Health & How is income associated with under-five mortality in 2022, and how
does the pattern differ by income group? & \file{materials/lab-4/health-brief.qmd} \\
Gender & Where is girls' secondary enrollment associated with lower
adolescent fertility? & \file{materials/lab-4/gender-brief.qmd} \\
Access & Which economies closed the combined electricity and internet access
gap fastest from 2000 to 2022? & \file{materials/lab-4/access-brief.qmd} \\
Climate & Which economies expanded renewables, reduced carbon intensity, and
ended with higher real income from 2000 to 2021? & \file{materials/lab-4/climate-brief.qmd} \\
\bottomrule
\end{tabularx}

\section*{Required brief}

Your rendered HTML should contain:
\begin{enumerate}
\item the question and a statement of the unit, sample, period, and source;
\item one table produced from the supplied CSV;
\item one labeled figure produced from the same analytical comparison;
\item one numerical finding written in your own words;
\item one limitation or inference the evidence does not support; and
\item an AI-use note describing the statement checked, independent R
evidence, and whether you accepted, revised, or rejected it.
\end{enumerate}

\section*{Route through the hour}

\begin{tabularx}{\textwidth}{@{}p{0.14\textwidth}p{0.24\textwidth}Y@{}}
\toprule
Minutes & Work & Result \\
\midrule
00--07 & Choose and plan & State the unit, variables, and function sequence. \\
07--20 & Construct table & One traceable briefing table. \\
20--33 & Construct figure & One labeled policy comparison. \\
33--40 & Interpret & One finding and one limitation in your words. \\
40--50 & Agent review & Codex reproduces the table and challenges one claim. \\
50--56 & Independent check & Test one material agent statement in R; decide
accept, revise, or reject. \\
56--60 & Render & A standalone HTML brief. \\
\bottomrule
\end{tabularx}

\section*{Before coding}

Write down the unit and period, the variables needed, the sequence of R
functions for the table, the figure mappings, one claim the evidence could
support, and one claim it could not support. The comments in the starter name
the relevant functions without completing the analysis.

\section*{Codex review}

Use Codex only after saving your table, figure, finding, and limitation. Replace
the bracketed filename in this request:

\begin{quote}\small
Read my saved [QMD filename] and the CSV it imports. Work read-only. Reproduce
my table from the stated sample; inspect the variables, units, and mappings in
my figure; then select the most consequential numerical or interpretive
statement in my Finding or Limitation. Cite the exact object and values
supporting your judgment. End with accept, revise, or reject and at most one
proposed revision. Do not edit files, fit a model, or make a causal claim.
\end{quote}

Choose one material statement from the response and test it with a small,
independent R calculation. In the Quarto file, record the agent's statement,
your check and result, and your final decision.

\section*{Interpretation reminders}

All four exercises are descriptive. The health and gender files are 2022
cross-sections. The access and climate files compare common endpoints. A
change between endpoints does not establish what happened in every
intervening year, and an observed association does not identify a policy
effect. Gross enrollment can exceed 100 by definition. The access index gives
equal weight to electricity and internet changes.

\section*{Completion}

Click \textbf{Render} in RStudio. Confirm that the HTML opens, includes the
table and figure, and contains no stale sentence that disagrees with the
output. There is no graded submission; keep the source QMD and rendered HTML as
a reusable record of your work.

\end{document}
